\chapter{Annexe C — Glossaire technique}

\begin{description}[leftmargin=3.5cm,style=nextline]

  \item[API REST]
    \textit{Application Programming Interface Representational State Transfer}.
    Interface de communication entre clients et serveur via des requêtes HTTP
    utilisant les méthodes GET, POST, PUT, PATCH et DELETE.

  \item[Bearer Token]
    Jeton d'accès transmis dans l'en-tête HTTP \texttt{Authorization: Bearer <token>}.
    Utilisé par Laravel Sanctum pour l'authentification stateless.

  \item[BFS]
    \textit{Breadth-First Search} (Parcours en largeur). Algorithme de parcours de
    graphe utilisé pour détecter les cycles dans le graphe de dépendances des tâches.

  \item[CPM]
    \textit{Critical Path Method} (Méthode du chemin critique). Algorithme de
    planification de projet identifiant la séquence de tâches dont le float est
    nul, déterminant ainsi la durée minimale du projet.

  \item[Drag-and-drop]
    Interaction UI permettant à l'utilisateur de glisser-déposer des éléments
    (cartes Kanban) d'une zone à une autre avec la souris.

  \item[Eager Loading]
    Stratégie Eloquent ORM chargeant en avance les relations associées
    (\texttt{with()}) pour éviter le problème des requêtes N+1.

  \item[Float (CPM)]
    Marge totale d'une tâche : différence entre sa date de début au plus tard (LST)
    et sa date de début au plus tôt (EST). Une tâche critique a un float de zéro.

  \item[GBM]
    \textit{Gradient Boosting Machine}. Algorithme d'ensemble ML qui construit
    séquentiellement des arbres de décision, chacun corrigeant les erreurs du
    précédent.

  \item[Kanban]
    Méthode de gestion visuelle du flux de travail, organisée en colonnes
    représentant les états d'avancement des tâches.

  \item[LLM]
    \textit{Large Language Model}. Modèle de langage de grande taille entraîné
    sur des corpus textuels massifs, capable de générer du texte en langage naturel.
    Ex. : Llama 3.3 70B.

  \item[Middleware]
    Composant logiciel Laravel interceptant les requêtes HTTP avant qu'elles
    n'atteignent le contrôleur. Ex. : \texttt{CheckRole}, \texttt{CheckProjectAccess}.

  \item[Optimistic Update]
    Technique UI consistant à mettre à jour l'interface immédiatement (avant la
    réponse serveur), puis à annuler en cas d'erreur.

  \item[RBAC]
    \textit{Role-Based Access Control}. Modèle de contrôle d'accès où les droits
    sont associés à des rôles, eux-mêmes attribués à des utilisateurs.

  \item[Random Forest]
    Algorithme d'ensemble ML créant un grand nombre d'arbres de décision sur des
    sous-ensembles aléatoires des données, puis agrégeant leurs prédictions.

  \item[Sanctum]
    Package Laravel fournissant une authentification légère par Bearer Token pour
    les SPA et les API REST.

  \item[SPA]
    \textit{Single Page Application}. Application web dont la navigation est
    entièrement gérée côté client (React Router), sans rechargement complet de page.

  \item[Sprint]
    Itération de longueur fixe (ici : 2 semaines) dans la méthode Scrum, à l'issue
    de laquelle un incrément fonctionnel livrable est produit.

  \item[Story Points]
    Unité relative d'estimation de l'effort dans Scrum, généralement exprimée en
    suite de Fibonacci (1, 2, 3, 5, 8, 13...).

  \item[VotingClassifier]
    Ensemble Scikit-learn combinant plusieurs classifieurs par vote pondéré
    (\textit{soft voting}) pour améliorer la robustesse des prédictions.

  \item[WebSocket]
    Protocole de communication bidirectionnel persistent entre client et serveur.
    Utilisé via Laravel Reverb pour les mises à jour en temps réel.

\end{description}
